\documentclass{hhuarticle}
% \usepackage[ngerman,english]{babel} % English
\usepackage[english,ngerman]{babel} % Deutsch
\usetikzlibrary{babel}
\usepackage{minted}

\title{Haskell - Angewandter \glqq abstract nonsense\grqq}

% Vorlange von Jannik Dunkelau und Joshua Schmid
\author{Luca Leon Happel} % keine Matrikelnummer eintragen

\subject{Bachelorseminar: Programmiersprachen} % Optional.
\semester{Wintersemester 2022}

% \printblackwhitetrue % Vereinfacht Schwarzweißdruck (Logo, Plots, ...).

%%%%%%%%%%%%%%%%%%%%%%%%%%%%%%%%%%%%%%%%%%%%%%%%%%%%%%%%%%%%%%%%%%%%%%%%%%%%%%%%
%% LaTeX Packages in Nutzung                                                  %%
%%                                                                            %%
%% Im folgenden können Sie für die Niederschrift benötigte LaTeX-Pakete       %%
%% einbinden.                                                                 %%
%% Diese Vorlage kommt bereits mit einigen nützlichen inkludierten Paketen.   %%
%%%%%%%%%%%%%%%%%%%%%%%%%%%%%%%%%%%%%%%%%%%%%%%%%%%%%%%%%%%%%%%%%%%%%%%%%%%%%%%%

%% Häufig benutzte mathematische Packages.
\usepackage{amsfonts}
\usepackage{amsmath}
\usepackage{amssymb}

\usepackage{enumitem} % Leichter konfigurierbare enumerate-Umgebungen.
\usepackage{subcaption} % Unterteilung von Figures in Subfigures.
\usepackage[colorlinks]{hyperref} % Klickbare Links (z.B. Inhaltsverzeichnis).
\sethyperrefhhucolors{} % Setzt den Farbsatz der HHU für hyperref.
\usepackage{url} % \url Kommando für Darstellung von Links
\usepackage{csquotes} % Improved quoting.

%% Tabellen
\usepackage{tabularx} % tabularx Umgebung für mehr Kontrolle über Tabellen.
\usepackage{booktabs} % \toprule, \midrule, \bottomrule
\usepackage{multirow}
\usepackage{multicol}
\usepackage{longtable} % Große Tabellen gehen über mehrere Seiten.

\usepackage{amsthm}
\usepackage{tikz-cd}
\usepackage{blindtext} %This package generates automatic text
\usepackage{epigraph}


\theoremstyle{definition}
\newtheorem{definition}{Definition}[section]

\theoremstyle{theorem}
\newtheorem{theorem}{Theorem}[section]

%% Quellcode
\usepackage{listings} % Einbindung von Code.
\setlstlistingstyle{} % Kosmetische Einstellungen
% Sprachabhängige Bezeichnung.
\iflanguage{ngerman}{\renewcommand{\lstlistingname}{Quellcode}}{}

%% Algorithmen in Pseudocode
\usepackage{algorithm} % Float-Umgebung für angegebene Algorithmen.
\usepackage{algorithmicx} % Angabe von Algorithmen in Pseudocode.
\usepackage{algpseudocode} % Standart Pseudocode-Elemente für Algorithmen.
\algsmallfont{}
% Sprachabhängige Bezeichnung.
\iflanguage{ngerman}{\floatname{algorithm}{Algorithmus}}{}

%% Intelligenteres Referenzieren mittels \cref.
\usepackage[capitalize,noabbrev,ngerman]{cleveref}
%\pagenumbering{gobble}

\definecolor{pblue}{rgb}{0.13,0.13,1}
\definecolor{pgreen}{rgb}{0,0.5,0}
\definecolor{pred}{rgb}{0.9,0,0}
\definecolor{pgrey}{rgb}{0.46,0.45,0.48}
\lstset{language=Java,
  showspaces=false,
  showtabs=false,
  breaklines=true,
  showstringspaces=false,
  breakatwhitespace=true,
  commentstyle=\color{pgreen},
  keywordstyle=\color{pblue},
  stringstyle=\color{pred},
  basicstyle=\ttfamily,
  moredelim=[il][\textcolor{pgrey}]{$$},
  moredelim=[is][\textcolor{pgrey}]{\%\%}{\%\%}
}

\begin{document}

  \maketitle

  \begin{abstract}
    % Was ist das Thema
    Haskell ist eine Sprache, welche vor mehr als 30 Jahren zuerst erschien
    und dennoch selbst heutzutage ihrer Zeit in manchen Bereichen voraus ist.
    % Warum ist das Thema interessant
    Trotz dessen ist sie in der Industrie wenig verbreitet, obwohl
    Programmierer viel von ihr lernen können.
    % Wie wurde es umgesetzt
    Inwiefern Programmierer davon profitieren können,
    wird in dieser Arbeit untersucht, indem
    Parallelen zwischen Haskell und anderen Sprachen gezogen werden.
    % Welche Ergebnisse wurden erzielt
    Wir erkennen dabei, dass Haskell mächtige Strukturen besitzt, welche
    in anderen Sprachen nicht oder nur teilweise existieren, da Haskell
    Datenstrukturen auf einem viel abstrakteren Level behandelt,
    wodurch Probleme mit Haskell weitaus effizienter und eleganter
    gelöst werden können.
    % Welche Gründe dies hat und warum sich dies ändern sollte, insbesondere
    % mit Blick auf Programmierer 
    % nicht nur Haskell, sondern
    % auch andere Sprachen, sowie Programmierer davon profitieren können,
    % wird in dieser Arbeit erläutert.
  \end{abstract}

\tableofcontents


  \section{Einleitung}

  Im September 1987 wurde die
  \glqq Conference on Functional Programming Languages and Computer Architecture\grqq 
  in Portland, Oregon, USA, abgehalten. Ziel war es, die zersplitterte
  funktionale Programmiersprachen-Community zu vereinen und eine
  gemeinsame Standardsprache zu entwickeln.
  Die Sprache sollte alle Vorteile der bis dahin existierenden
  funktionalen Sprachen, insbesondere der zu der Zeit populären Sprache
  \glqq Miranda\grqq , vereinen und gleichzeitig die Nachteile beseitigen.
  Die Sprache, welche dabei entstand, ist Haskell \cite{haskellHistory}.

  Während ihrer Entwicklung wurde Haskell stark von Akademikern beeinflusst,
  wobei insbesondere Themen wie Typentheorie, Kategorientheorie und das Lambda-Kalkül
  eine große Rolle spielten. Daraus resultierte, dass Haskell
  einige Besonderheiten aufweist, welche es von anderen Sprachen unterscheidet.
  Diese Besonderheiten sind es, welche Haskell zu einer so interessanten
  Sprache machen, da somit Haskell eine der wenigen reinen
  Programmiersprachen ist, also keine Seiteneffekte vorhanden sind.
  Eine unintuitive Konsequenz daraus ist, dass Haskell keine Variablen besitzt,
  welche verändert werden können, da es keine Zuweisung gibt,
  da Zuweisungen in Haskell Seiteneffekte erzeugen würden. Somit
  existieren in Haskell nur Konstanten, welche nicht verändert werden können.
  Dies wiederum führt dazu, dass klassische Schleifen in Haskell nicht existieren,
  da diese durch Seiteneffekte, insbesondere dem Inkrementieren einer
  Laufvariable, implementiert werden. Stattdessen werden in Haskell
  Funktionen verwendet, welche rekursiv aufgerufen werden, um
  die gewünschte Schleifenfunktionalität zu erreichen.

  Statt der Objektorientierung verwendet Haskell die funktionale
  Programmierung, was durch die Reinheit bedeutet, dass Haskell
  sehr gut durch kategorientheoretische Konzepte betrachtbar ist.
  \grqq Kategorientheoretische Konzepte\grqq  beziehen sich auf die
  Kategorientheorie, welche eine Theorie ist, welche
  eine alternative Grundlage für die Mathematik darstellt, im Gegensatz
  zu der klassischen Mengenlehre.
  Oft werden Konzepte aus der Kategorientheorie auch als \glqq abstract nonsense \grqq
  bezeichnet \cite{henderson2008generalized}.
  Diese besagt, dass es Kategorien gibt, in welchen es sogenannte Objekt
  gibt, welche durch Morphismen verbunden sind. Existieren
  Objekte $A$, $B$ und $C$ in einer Kategorie $\mathcal{C}$,
  sowie Morphismen $f: A \to B$ und $g: B \to C$, so gibt es einen
  Morphismus $g \circ f: A \to C$ und außerdem muss es für jedes
  Objekt $A$ einen Morphismus $id_A: A \to A$ geben, welcher
  das Objekt $A$ nicht verändert. In dem Fall von Haskell sind die
  Objekte die Typen und die Morphismen die Funktionen, wobei dies noch
  weitreichende Konsequenzen hat, welche im Folgenden noch aufgefasst werden.
  Funktionale Programmierung bedeutet, dass Programme durch Funktionen
  definiert werden, welche Daten als Eingabe erhalten und Daten als
  Ausgabe liefern. Geschickte Komposition von Funktionen ermöglicht
  es, komplexe Probleme durch hintereinander aufrufende einfache
  Funktionen zu lösen.
  Dies hat zur Folge, dass Programmierer, welche Haskell lernen,
  automatisch besser in abstrakter Mathematik werden. Außer der Reinheit
  und der funktionalen Programmierung, ist Haskell eine
  polymorphe, statisch-typisierte Programmiersprache mit
  lazy-evaluation, wobei diese Begriffe im Folgenden erläutert werden.

  Polymorphie bedeutet, dass Funktionen und Datenstrukturen
  generisch sind, was bedeutet, dass sie mit verschiedenen Typen verwendet werden können.
  Statische Typisierung bedeutet, dass der Typ einer Variable
  zur Kompilierzeit festgelegt wird, wodurch Fehler, welche durch
  dynamische Typisierung entstehen, vermieden werden. Durch das
  algebraische Typensystem von Haskell, was bedeutet, dass Typen $A$ und
  $B$ zu Produkttypen $A \times B$ und Summentypen $A | B$ kombiniert werden können,
  ist es möglich, Funktionen und Datenstrukturen mit verschiedenen
  Typen zu kombinieren. Hier sehen wir wieder, dass ein starker Fokus
  auf dem Kombinieren von einfacheren Strukturen liegt, wodurch komplexere
  Strukturen entstehen. Da Haskell auch eine funktionale Programmiersprache ist,
  haben auch Funktionen einen Typen, wodurch Funktionen, Konstanten und Variablen
  alle durch die starke Typisierung von Haskell definiert werden können
  und somit eine Typenstruktur aufweisen. Da wir bereits wissen, dass
  Haskell eine stark-typisierte Programmiersprache ist, bedeutet dies,
  dass Programmierfehler in Haskell fast immer von dem Compiler erkannt werden.
  Sollte eine Funktion mit einem falschen Typen aufgerufen werden, oder
  einen falschen Typen zurückgeben, so wird dies vom Compiler erkannt.
  Diese Eigenschaft von Haskell ist sehr nützlich, da es so möglich ist,
  sicherheitskritische Programme zu schreiben, welche mit erhöhter
  Wahrscheinlichkeit korrekt sind als Programme in anderen Sprachen.

  Die lazy-evaluation von Haskell bedeutet, dass Ausdrücke erst dann
  ausgewertet werden, wenn sie benötigt werden. Dies hat zur Folge,
  dass Ausdrücke, welche nie benötigt werden, nicht ausgewertet werden.
  Dies ist sehr nützlich, da es so möglich ist, unendliche Datenstrukturen
  zu definieren, welche nur dann ausgewertet werden, wenn sie benötigt werden.
  Beispielsweise kann eine Liste von allen Primzahlen definiert werden,
  wobei man nun die $n$-te Primzahl aus dieser Liste erhalten kann,
  ohne dass die Liste vollständig ausgewertet werden muss.
  Auch wird Code somit idiomatischer, da man analog zur natürlichen
  Sprache \glqq die $n$-te Primzahl\grqq  verwenden kann, statt diese
  durch eine Funktion zu erhalten, welche die $n$-te Primzahl zurückgibt.

  Diese Eigenschaften von Haskell machen es zu einer sehr ungewöhnlichen
  Programmiersprache, welche sich von anderen Programmiersprachen
  wie C, Java oder Python unterscheidet. Diese Unterschiede werden
  im Folgenden noch näher erläutert, wobei wir sehen werden,
  dass viele Konzepte von Haskell über die Jahre durch ihre praktische
  Anwendung in anderen Programmiersprachen übernommen wurden.
  Der tiefe mathematische Hintergrund von Haskell ist jedoch
  immer noch außergewöhnlich und erlaubt durch die abstrakten
  Strukturen darin, Probleme sehr elegant zu lösen, wobei dies auch
  ein zweiseitiges Schwert ist, da es für Programmierer, welche
  nicht darin geschult sind, sehr schwer zu verstehen ist. Hier werden
  wir auch sehen, dass andere Programmiersprachen, welche die
  Konzepte von Haskell übernommen haben, diese Konzepte nicht
  so abstrakt und rigoros implementiert haben wie Haskell, wodurch
  die Konzepte dort auf bestimmte Anwendungen spezialisiert sind,
  und in ihrem Umfang eingeschränkt sind.
  Als Beispiel dafür werden wir den Typen \verb|Option<u32>| in
  Rust betrachten, welcher ähnlich dem Typen \verb|Maybe Int| in Haskell
  ist, jedoch aufgrund seiner fehlenden Funktor-Struktur unhandlicher
  ist und mehr Fallunterscheidungen erfordert.

  % Diese Eigenschaften machen Haskell zu einer sehr mächtigen Sprache, welche
  % es ermöglicht, elegante und effiziente Lösungen für viele Probleme zu finden.
  % Durch ihren starken mathematisch, theoretischen Grundaufbau besitzt Haskell
  % Strukturen tief in sich verankert, welche noch heutzutage
  % von anderen Programmiersprachen teilweise übernommen werden.
  % Wie genau diese Strukturen in Haskell implementiert sind, wird im Folgenden
  % erläutert. Außerdem werden wir sehen, wie Spezialisierungen 


  % Jede Ausarbeitung beginnt mit einer guten Einleitung.
  % Hier liegt der Fokus auf den Grundlagen,
  % sodass klar ist, worum es geht, was das Thema behandelt und
  % wieso dies in erster Linie interessant ist.
  % Detaillierte Informationen folgen in den späteren Abschnitten.

  % Diese grundlegende Einleitung erlaubt auch,
  % dass spätere Informationen hierauf aufbauen können.
  % Dadurch lässt sich verhindern,
  % dass zwei unterschiedliche Abschnitte die gleichen Grundlagen definieren
  % und somit im Dokument doppeln.


  \section{Grundlagen von Haskell}
  
  Wir möchten Haskell mit anderen Programmiersprachen vergleichen.
  Interessant ist, dass wir Konzepte, welche in Haskell verwendet werden
  und aus der Kategorientheorie stammen, dazu verwenden können, um dieses
  Ziel zu verwirklichen.

  \subsection{Typen und Typklassen}

  Betrachten wir zuerst eine Typenklasse, welche in Haskell definiert ist.
  Dazu müssen wir jedoch zuerst definieren, was eine Typenklasse ist.

  \begin{definition}
    Eine Typenklasse ist eine Sammlung von Typen, welche eine bestimmte
    Eigenschaft besitzen.
  \end{definition}

  Nach der Einleitung wissen wir bereits, was ein Typ ist. Diese kommen
  aus der Typentheorie und ordnen Werten in unserer Programmiersprache
  eine Eigenschaft genannt \glqq Typ\grqq  zu. Allgemeiner bezeichnen wir als
  ein \glqq Typensystem\grqq  eine konsistente Theorie $T$, für welche zu jeder Aussage $\phi$
  der Form \glqq $x$ hat den Typ $y$\grqq , mit $x$ ein Wert und $y$ ein Typ,
  entweder $\phi$ gilt in $T$, oder $\neg \phi$ gilt in $T$, wobei \glqq $\neg$\grqq 
  die logische Negierung ist, gilt. Dabei bedeutet \glqq konsistente\grqq , dass
  es keine Kontradiktion in dieser Theorie gibt, also keine Aussage $\phi$
  existiert, für welche sowohl $\phi$ in $T$ als auch in $\neg \phi$ in $T$ gilt.
  Ein Beispiel für einen Typ ist der Typ $Bool$ welcher
  den Werten $True$ und $False$ in Haskell zugeordnet ist. Ein weiterer Typ wäre
  der Typ $Int$, welcher den ganzen Zahlen zugeordnet ist.
  Ein Beispiel für eine Typentheorie ist die Typentheorie von Haskell,
  welche unter anderem die Typen $Bool$ und $Int$ enthält. Inkonsistent,
  also nicht konsistent, wäre die Typentheorie von Haskell, wenn
  beispielsweise die Aussage \glqq $True$ hat den Typ $Int$\grqq  in dieser Theorie
  und die Aussage \glqq $True$ hat nicht den Typ $Int$\grqq  gelten würde, was
  jedoch nicht gilt.

  Die Definition einer Typenklasse als Sammlung von Typen ist jedoch auf
  den ersten Blick nicht sonderlich aufschlussreich. Der Begriff \glqq Sammlung\grqq 
  bezieht sich in dem Kontext von Haskell auf eine Menge von Typen im
  Sinne der Mengenlehre. Wir können sogar weiter eingrenzen und stark davon
  ausgehen, dass die Menge von Typen, welche eine Typenklasse definiert,
  endlich in Haskell ist, da ein Programm in Haskell nur endlich viele
  Typen definieren kann.

  Wir können nun die Implementation von Typen und Typenklassen in Haskell
  genauer betrachten. Wir können in \cref{fig:mytype} sehen, wie ein
  einfacher Typ in Haskell definiert wird und in \cref{fig:myfunction} eine
  Funktion, welche einen Typen als Parameter verwendet. Insbesondere
  bezeichnen wir die erste Zeile \verb|f :: MyType -> Bool| als \glqq Typsignatur\grqq 
  und die zweite Zeile \verb|f (MyConstruktor b i) = b| als \glqq Definition\grqq  der Funktion.
  
  \begin{figure}[h]
    \begin{minted}[frame=single]{haskell}
    data MyType = MyConstruktor Bool Int
    \end{minted}
    \caption{Definition eines Typens genannt \glqq MyType\grqq  mit Konstruktor \glqq MyConstruktor\grqq , welcher einen \glqq Bool\grqq  und einen \glqq Int\grqq  speichert \cite{constructorHaskellWiki}}%
    \label{fig:mytype}
  \end{figure}

  \begin{figure}[h]
    \begin{minted}[frame=single]{haskell}
    f :: MyType -> Bool
    f (MyConstruktor b i) = b
    \end{minted}
    \caption{Definition einer Abbildung, welche einen Wert von Typ \glqq MyType\grqq  auf einen Wert von Typ \glqq Bool\grqq  schickt.}%
    \label{fig:myfunction}
  \end{figure}

  Eine Typenklasse wird nun wie in \cref{fig:mytypeclass} definiert. Wir
  können sehen, dass eine Typenklasse eine Sammlung von Typen
  parametrisiert durch einen Typenvariablen $a$ ist, wobei
  eine Funktion \verb|myFunction :: a -> Bool| existieren muss.
  Diese Typenklasse ist jedoch bisher nur eine Definition, wobei kein
  Typ ein Teil dieser Typenklasse ist. Damit wir einen Typen in diese
  Typenklasse hinzufügen können, müssen wir beweisen, dass ein gegebener
  Typ die Eigenschaften der Typenklasse erfüllt. Dies wird in \cref{fig:mytypeinstance}
  gezeigt. Wir können sehen, dass wir die Typenklasse \verb|MyTypeClass| mit dem Typen
  \verb|MyType| verbinden, indem wir den Typen \verb|MyType| als
  Instanz der Typenklasse \verb|MyTypeClass| definieren.
  Insbesondere haben wir für alle Typensignaturen, welche nach dem \verb|where|
  in \cref{fig:mytypeclass} stehen, eine Definition, welcher der
  entsprechenden Typensignatur zustimmt, in der Instanziierung gegeben.

  \begin{figure}[h]
    \begin{minted}[frame=single]{haskell}
    class MyTypeClass a where
      myFunction :: a -> Bool
    \end{minted}
    \caption{Definition einer Typenklasse, welche einen Typen \glqq a\grqq  als Parameter besitzt.}%
    \label{fig:mytypeclass}
  \end{figure}

  \begin{figure}[h]
    \begin{minted}[frame=single]{haskell}
    instance MyTypeClass MyType where
      myFunction (MyConstruktor b i) = b
    \end{minted}
    \caption{Definition einer Typenklasse, welche einen Typen \glqq a\grqq  als Parameter besitzt.}%
    \label{fig:mytypeinstance}
  \end{figure}

  Als ein pathologisches Beispiel für eine Typenklasse, welche von jedem
  Typen $a$ implementiert werden kann, wobei \glqq implementiert\grqq  gleichzusetzen
  mit \glqq $a$ ist ein Typ in der entsprechenden Typenklasse\grqq  ist, ist die
  Typenklasse \verb|Trivial|, welche in \cref{fig:trivialTypeClass} definiert ist.
  Die entsprechende Implementation für den Typen \verb|MyType| ist in \cref{fig:trivialTypeClassInstance},
  wobei wir sehen, dass wir keine Implementationen
  für Funktionen vornehmen müssen.

  \begin{figure}[h]
    \begin{minted}[frame=single]{haskell}
    class Trivial a
    \end{minted}
    \caption{Definition einer pathologischen Typenklasse, welche für jeden Typen \glqq a\grqq  als Parameter ohne Einschränkung implementiert werden kann.}%
    \label{fig:trivialTypeClass}
  \end{figure}

  \begin{figure}[h]
    \begin{minted}[frame=single]{haskell}
    instance Trivial MyType
    \end{minted}
    \caption{Definition einer Typenklasse, welche einen Typen \glqq a\grqq  als Parameter besitzt.}%
    \label{fig:trivialTypeClassInstance}
  \end{figure}

  \subsection{Funktoren}

  Als besondere Typenklasse ist die \verb|Functor| Typenklasse \cite{functorHaskellWiki} zu
  erwähnen. Die \verb|Functor| Typenklasse ist eine Typenklasse, welche
  im \glqq Prelude\grqq  von Haskell
  \footnote{der Prelude ist gewissermaßen die Standartbibliothek von
  Haskell, welche automatisch in jedem Haskellprogramm eingebunden wird}
  definiert ist. Diese Definition können wir in
  \cref{fig:functorTypeClass} genauer betrachten.
  Die erste Zeile ist ein \glqq Type-Constraint\grqq  und besagt, dass
  nur Typen der Art \verb|*->*| in die \verb|Functor| Typenklasse
  implementiert werden können. Die \glqq Art\grqq  ist ein Konzept
  aus der Typentheorie \cite{kindHaskellWiki}, welches wir hier nicht näher betrachten.
  Das einzige wichtige für uns ist, dass die Art \verb|*->*|
  besagt, dass der Typenparameter $a$ ein Typenkonstruktor ist,
  welcher einen Typen reinnimmt und einen Typen zurückgibt.
  Hierbei sind Typenkonstruktoren bereits aus \cref{fig:mytype} bekannt,
  wobei dort der Typenkonstruktor \verb|MyConstruktor| war. Dieser
  hätte die Art \verb|*|, da er keine freien Typenparameter gibt,
  wohingegen \verb|data MyType a b = MyConstruktor a b| von der
  Art \verb|*->*->*| ist, da er zwei freie Typenparameter hat.

  \begin{figure}[h]
    \begin{minted}[frame=single]{haskell}
    type Functor :: (* -> *) -> Constraint
    class Functor f where
      fmap :: (a -> b) -> f a -> f b
      (<$) :: a -> f b -> f a
      {-# MINIMAL fmap #-}
        -- Defined in 'GHC.Base'
    \end{minted}
    \caption{Definition der \glqq Functor\grqq  Typenklasse.}%
    \label{fig:functorTypeClass}
  \end{figure}

  In der Definition der \verb|Functor| Typenklasse sehen wir, dass
  nach dem \verb|where| eine Typensignatur \verb|fmap| und eine
  Typensignatur \verb|<$| definiert sind. Diese Typensignaturen
  sind die Funktionen, welche die Instanzen der \verb|Functor| Typenklasse
  implementieren müssen.
  Die Typensignatur \verb|fmap| besagt, dass
  wir eine Funktion \verb|fmap| definieren müssen, welche eine Funktion
  vom Typen \verb|(a -> b)| und einen Wert vom Typen \verb|f a| einnimmt
  und einen Wert vom Typen \verb|f b| zurückgibt. Wichtig für das
  Verständnis dieser Typensignaturen, welche mehr als einen Pfeil besitzen,
  ist das sogenannte Konzept des Curryings
  \footnote{Auch \glqq Schönfinkeln\grqq  in dem Bereich der Linguistik gennant}.

  \subsubsection{Currying}

  Das Konzept des Curryings ist ein Konzept aus der Mathematik, welches
  in der Programmierung verwendet wird, um Funktionen mit mehreren
  Argumenten mittels einer natürlichen Bijektion zwischen $\text{Hom}$-Mengen
  elegant zu definieren. Was wir damit meinen, möchten wir in diesem
  Unterabschnitt näher erläutern, genauso wie, welche Konsequenzen
  das Curryings für die Typensignaturen der \verb|Functor| Typenklasse
  hat.

  Nach der Einleitung dieser Arbeit wissen wir bereits, was eine
  Kategorie $\mathcal{C}$ ist. Wir wissen, dass diese Kategorie
  Objekte $A, B$ beinhalten kann, sowie Morphismen $f:A\rightarrow B$.
  Wir einigen uns auf die Konvention, dass wir $A\in \text{ob}(\mathcal{C})$
  für ein Objekt $A$ in der Kategorie $\mathcal{C}$ schreiben und
  $f:A\rightarrow B \in \text{Hom}_\mathcal{C}(A, B)$ für
  einen Morphismus von $A$ nach $B$ in der Kategorie $\mathcal{C}$.
  Allgemein sind $\text{ob}(\mathcal{C})$ und $\text{Hom}_\mathcal{C}(A, B)$
  echte Klassen, also keine Mengen im Sinne der Mathematik.
  Wir werden jedoch nicht näher auf den Unterschied zwischen
  Mengen und Klassen eingehen, da dies nicht relevant für die
  Kategorie \textbf{Hask} ist, welche die Kategorie der Typen und 
  Funktionen in Haskell ist\footnote{Per Konvention nehmen wir an,
  dass Funktionen in Haskell terminieren, also endliche Laufzeit haben}.
  Die genaue Konstruktion der Kategorie \textbf{Hask}
  lässt sich in der offiziellen Dokumentation von Haskell
  \footnote{\path{https://wiki.haskell.org/Hask#\%22Platonic\%22_Hask}}
  nachlesen. Wichtig ist nur, dass Typen in Haskell Objekte in der
  Kategorie \textbf{Hask} sind und Funktionen zu Morphismen in der
  Kategorie \textbf{Hask} unter bestimmten Voraussetzungen korrespondieren.
  In der Kategorie \textbf{Hask} gilt, dass $\text{ob}(\textbf{Hask})$
  und $\text{Hom}_\textbf{Hask}(A, B)$ Mengen für alle $A, B \in \text{ob}(\textbf{Hask})$ sind.

  \begin{definition}
    Wir definieren den Exponenten $Z^Y$ von Objekten $Z, Y\in \text{ob}(\textbf{Hask})$
    als die Menge aller Morphismen $h: Z \to Y$ vom Objekt $Z$ nach $Y$.
  \end{definition}

  Wir können nun das Konzept des Curryings in der Kategorie \textbf{Hask}
  definieren. Dazu betrachten wir uns folgende Eigenschaft der Kategorie
  \textbf{Hask}:

  \begin{theorem}
    Für alle Objekte $X, Y, Z \in \text{ob}(\textbf{Hask})$ gilt:
    \begin{equation}
      \text{Hom}_\textbf{Hask}(X \times Y, Z) \cong \text{Hom}_\textbf{Hask}(X, Y^Z)
    \end{equation}
  \end{theorem}

  \begin{proof}
    Sei $f: X \times Y \to Z$ ein Morphismus in der Kategorie \textbf{Hask}.
    Dieser korrespondiert zu einer Funktion $f: X \times Y \to Z$ in Haskell,
    also eine Funktion, welche ein Tupel aus $X$ und $Y$ einnimmt und etwas
    vom Typ $Z$ zurückgibt. Wir definieren nun eine Funktion
    $g: X \to Y^Z$ als $g(x) = f(x, \cdot)$, wobei $\cdot$ ein Platzhalter
    für ein Argument vom Typ $Y$ ist, also
    $f(x, \cdot) = (y \mapsto f(x, y))$.

    Es ist klar, dass $g$ eine Funktion vom Typ $X \to Y^Z$ ist, da
    $g$ einen Wert vom Typ $X$ einnimmt und eine Funktion von $Y$ nach $Z$
    zurückgibt. Dies sind aber genau die Werte vom Typ $Y^Z$.

    Allgemein können wir eine Abbildung
    $\text{Hom}_\textbf{Hask}(X \times Y, Z) \tilde\to \text{Hom}_\textbf{Hask}(X, Y^Z)$
    konstruieren, indem wir für jeden Morphismus $f: X \times Y \to Z \in \text{Hom}_\textbf{Hask}(X \times Y, Z)$
    die Abbildung $\phi(f) = g$ definieren, wobei $g$ die Funktion
    $g: X \to Y^Z, \> g(x) = (y \mapsto f(x,y)$ ist. Die Bijektivität
    ist nach Konstruktion gegeben, was wir durch $\tilde \to$ symbolisieren.
  \end{proof}

  Wir können noch viel mehr über diese Abbildung $\phi$ sagen, jedoch
  beschränken wir uns darauf, dass dies als Implikation bedeutet, dass
  die Funktionen in \cref{eq:currying1} praktisch identisch sind.
  Für den interessierten Leser bietet sich hierbei \cite[S. 53, Proposition 3.3.2]{CATEGORIES_TYPES_AND_STRUCTURES} an.
  Man kann also \verb|f| und \verb|g| als die gleiche Funktion
  betrachten, da man durch Anwenden von $\phi$ auf \verb|f| die Funktion
  \verb|g| erhält und umgekehrt durch Anwenden von $\phi^{-1}$ auf \verb|g|
  die Funktion \verb|f| erhält. Wichtig ist hierbei,
  dass wir \verb|g x|, also die Funktion \verb|g| angewandt auf
  den Parameter \verb|x|, durch eine Lambda-Funktion
  \verb|(\y -> f (x y))| definiert haben, welche also ein \verb|y| nimmt
  und \verb|f| auf \verb|(x, y)| angewandt ausgibt.

  \begin{figure}[h]
    \begin{minted}[frame=single]{haskell}
    f :: (a, b) -> c
    g :: a -> (b -> c)
    g x = \y -> f (x, y)
    \end{minted}
    \caption{Typensignaturen von Funktionen, welche äquivalent bezüglich curryings sind}%
    \label{fig:currying1}
  \end{figure}

  Wir erkennen aus dieser Tatsache, dass es effizienter ist, die Funktionen
  \verb|f| und \verb|g| direkt ohne Klammern zu schreiben. So erhalten wir
  eine Kurzschreibweise für die Funktionen \verb|f| und \verb|g|, welche
  wir in \cref{fig:curryingf} als Funktion \verb|h| sehen können.
  Dem Verständnis geschuldet, haben wir auch eine äquivalente Definition
  von \verb|h| in \cref{fig:curryingg} gegeben.

  \begin{figure}[h]
    \begin{minted}[frame=single]{haskell}
    f :: (a, b) -> c
    g :: a -> (b -> c)
    g x = \y -> f (x, y)

    h :: a -> b -> c
    h x y = f (x, y)
    \end{minted}
    \caption{Kurzschreibweise resultierend aus dem Currying basierend basierend auf der Definition von f}%
    \label{fig:curryingf}
  \end{figure}

  \begin{figure}[h]
    \begin{minted}[frame=single]{haskell}
    f :: (a, b) -> c
    g :: a -> (b -> c)
    g x = \y -> f (x, y)

    h :: a -> b -> c
    h x y = (g x) y
    \end{minted}
    \caption{Kurzschreibweise resultierend aus dem Currying basierend basierend auf der Definition von g}%
    \label{fig:curryingg}
  \end{figure}

  Die Vorteile hierbei sind vielfältig. Ein offensichtlicher Vorteil ist,
  dass Funktionen effizienter und lesbarer dargestellt werden können,
  wie wir in \cref{fig:curryingExample} sehen können. Hier nehmen wir
  ans, dass die Funktion \verb|callServer| zwei Parameter vom Typ \verb|String| nimmt
  und einen Wert vom Typ \verb|String| zurückgibt. Sie könnte beispielsweise
  eine URL auf einen Server, definiert im ersten Parameter, mit
  einem bestimmten Header, definiert im zweiten Parameter, aufrufen und
  den Inhalt der Seite zurückgeben.

  \begin{figure}[h]

    \begin{minted}[frame=single]{haskell}
    callServer :: String -> String -> String
    callServer url header = ...

    apiCall :: String -> String
    apiCall = callServer "https://api.example.com"

    apiCallWithHeader :: String
    apiCallWithHeader = apiCall "example header"
    \end{minted}
    \caption{So binden Sie Code ein}%
    \label{fig:curryingExample}
  \end{figure}

  Des Weiteren kann man durch currying Funktionen höhere Ordnung
  eleganter verwenden. Am Beispiel der Funktion \verb|map : (a -> b) -> [a] -> [b]|,
  welche eine Funktion vom Typ \verb|a -> b| auf eine Liste vom Typ \verb|[a]| anwendet
  und eine Liste vom Typ \verb|[b]| zurückgibt, können wir sehen, dass
  das Currying in Anbetracht der Funktion \verb|+ : Num a => a -> a -> a| sehr nützlich ist.
  Hierbei ist \verb|+| eine Funktion vom Typ \verb|Num a => a -> a -> a|,
  was bedeutet, dass \verb|+| zwei Werte des Typs \verb|a| nimmt und einen Wert
  des Typs \verb|a| zurückgibt, wobei \verb|a| ein Typ ist, der die Klasse \verb|Num| implementiert.
  \verb|Num| ist eine Typenklasse, welche die Grundrechenarten implementiert.
  Also sagen wir mit \verb|Num a => ...| aus, dass \verb|+| auf dem
  Typen \verb|a| definiert sein muss.

  \begin{figure}[h]

    \begin{minted}[frame=single]{haskell}
    increaseByOne :: [Int] -> [Int]
    increaseByOne = map (+ 1)
    \end{minted}
    \caption{Wir komponieren die map Funktion geschickt mit einer partiellen Anwendung der gecurryten Funktion + um eine Funktion zu erhalten, welche alle Elemente einer Liste von Zahlen um 1 zu erhöhen.}%
    \label{fig:currying1}
  \end{figure}

  \subsubsection{Anwendung der Functor Typenklasse}

  Nachdem wir nun das notwendige Wissen über currying haben um 
  die Typensignaturen der Functor Typenklasse in \cref{fig:functorTypeClass}
  zu verstehen, können wir uns nun mit der Implementierung der Functor Typenklasse
  beschäftigen und herausfinden, warum diese so wichtig ist.

  \begin{figure}[h]
    \begin{minted}[frame=single]{haskell}
    fmap (\x -> x+1) [1..10]
    \end{minted}
    \caption{Anwendung der kanonischen Funktor Instanz auf der Liste von ganzen Zahlen $z$ mit $0\le z < 10$, wobei jede Zahl in der Liste um eins erhöht wird.}%
    \label{fig:fmapExample}
  \end{figure}

  Betrachten wir dafür \cref{fig:fmapExample}, wobei wir erkennen,
  dass wir durch \verb|fmap| eine Funktion vom Typ \verb|a -> b| auf
  eine Liste vom Typ \verb|[a]| anwenden können. Wir haben also unsere
  Funktion vom Typ \verb|a -> b| auf einen Typen, welcher unseren
  ursprünglichen Typen \verb|a| beinhaltet \glqq hochgehoben\grqq .
  Das wichtige hierbei ist, dass wir erkannt haben, dass Listen von einem
  beliebigen Typen \verb|[a]| eine \glqq kanonische\grqq , beziehungsweise offensichtliche,
  Möglichkeit besitzen, Funktionen vom Typ \verb|a -> b| auf sich anzuwenden.

  Dieses Konzept tritt überall in der Programmierung auf. Wir möchten
  mit Strukturen arbeiten, indem wir ihre Bestandteile manipulieren.
  Um dieses Konzept zu festigen, betrachten wir einen weiteren Typen in
  Haskell, welcher durch seine Nützlichkeit auch in der Programmiersprache
  Rust anzutreffen ist. Dieser Typ wird in Haskell \verb|Maybe| genannt
  und in Rust \verb|Option| \cite{rustOption}.

  \begin{figure}[h]
    \begin{minted}[frame=single]{haskell}
      data Maybe a = Nothing | Just a

      meineDivision :: Int -> Int -> Maybe Int
      meineDivision x y = if y == 0 then Nothing else Just (x `div` y)
    \end{minted}
    \caption{Der Maybe Datentyp und wie man eine einfachere Funktion mit diesem schreibt.}%
    \label{fig:maybeType}
  \end{figure}

  Der Maybe Typ, wie er in \cref{fig:maybeType} definiert ist, ist auch
  Bestandteil des \verb|Prelude| Pakets und somit auch in jedem Haskell
  Programm verfügbar. Er ist ein Typ, welcher entweder den Wert \verb|Nothing|
  oder den Wert \verb|Just a| beinhaltet, wobei \verb|a| ein beliebiger Typ ist.
  Diese Definition basiert auf dem algebraischen Typensystem von Haskell,
  welches wir in der Einleitung kennengelernt haben.
  Außerdem lernen wir nun die \verb|if| Anweisung kennen, welche
  in Haskell durch \verb|if ... then ... else ...| verwendet wird.
  Diese ist klar verständlich und ähnlich wie in anderen Programmiersprachen.
  Die Funktion \verb|meineDivision| nimmt zwei Zahlen vom Typ \verb|Int| und
  gibt entweder \verb|Nothing| oder \verb|Just Int| zurück,
  je nachdem ob der zweite Parameter \verb|y| gleich $0$ ist oder nicht.
  Dies erfasst die allgemein bekannte Idee, dass \glqq etwas schlimmes passiert,
  wenn man durch $0$ teilt\grqq .

  \begin{figure}[h]
    \begin{minted}[frame=single]{haskell}
    erhoeheMaybe :: Maybe Int -> Maybe Int
    erhoeheMaybe Nothing = Nothing
    erhoeheMaybe (Just x) = Just (x + 1)
    \end{minted}
    \caption{Naive Definition von \glqq erhoeheMaybe\grqq }%
    \label{fig:maybeTypeAnwenden1}
  \end{figure}

  Wir können nun die Funktion \verb|erhoeheMaybe| wie in \cref{fig:maybeTypeAnwenden1} definieren, welche
  eine Funktion vom Typ \verb|Maybe Int -> Maybe Int| ist. Diese
  Funktion erhöht den Wert eines Wertes von Typ \verb|Maybe Int| um eins, falls
  dieser nicht \verb|Nothing| ist, also \verb|Just Int| ist und somit
  einen Wert vom Typ \verb|Int| beinhaltet.
  Diese Funktion lässt somit \verb|+ 1| auf kanonische Weise
  auf dem Typen \verb|Maybe Int| wirken, wobei dies ein Stichwort
  dafür ist, dass \verb|Maybe| eine Instanz der Functor
  Typenklasse ist. Sieht man in der Haskell-Dokumentation nach, so findet
  man heraus, dass \verb|Maybe| eine Instanz der Functor Typenklasse ist,
  und somit unser Instinkt richtig war, wodurch wir Funktionen wie \verb|(+ 1)|
  vom Typ \verb|Int -> Int| auf \verb|Maybe Int| anwenden können.
  Durch ein geschicktes Refactoring können wir die Funktion
  \verb|erhoeheMaybe| somit in eine Zeile umschreiben, statt eine
  Fallunterscheidung zu schreiben. Wir können dies in \cref{fig:maybeTypeAnwenden2}
  sehen.

  \begin{figure}[h]
    \begin{minted}[frame=single]{haskell}
    erhoeheMaybe = fmap (+ 1)
    \end{minted}
    \caption{Durchdachtere Definition von \glqq erhoeheMaybe\grqq }%
    \label{fig:maybeTypeAnwenden2}
  \end{figure}

  Das interessante ist, dass wir hierbei durch unser Refactoring nicht
  nur die Lesbarkeit verbessert haben, sondern auch deutlich weniger
  Code geschrieben haben, sowie durch das Erkennen der unterliegenden
  Struktur in der Lage sind, die Funktion \verb|erhoeheMaybe| statt
  nur auf \verb|Maybe Int| zu definieren, allgemeiner zu definieren.
  Indem wir die Funktion \verb|fmap| verwenden und somit die Funktorialität
  von \verb|Maybe| ausnutzen, mussten wir gar keinen konkreten Bezug auf
  den Typen Maybe mehr herstellen. Wir können also eine neue Typensignatur
  für \verb|erhoeheMaybe| definieren, welche in \cref{fig:maybeTypeAnwenden3}
  zu sehen ist.

  \begin{figure}[h]
    \begin{minted}[frame=single]{haskell}
    erhoeheMaybe :: (Functor f, Num b) => f b -> f b
    erhoeheMaybe = fmap (+ 1)
    \end{minted}
    \caption{Typensignatur von \glqq erhoeheMaybe\grqq  nach einem Refactoring. \glqq erhoeheMaybe\grqq  ist nun auf allen \glqq f\grqq  anwendbar, welche Funktoren sind und einen Typen \glqq b\grqq , welcher eine Zahl ist, im ersten Parameter haben. }%
    \label{fig:maybeTypeAnwenden3}
  \end{figure}

  Hiermit haben wir eine enorme Verallgemeinerung erreicht, wodurch unsere
  Funktion \verb|erhoeheMaybe| beispielsweise im Spezialfall von \verb|f| = \verb|[]|
  und \verb|b| = \verb|Int|, also auf Listen von Integern, angewendet werden kann.
  In diesem Spezialfall erhöht \verb|erhoeheMaybe| die Werte in einer Liste von Zahlen
  jeweils um eins, was bedeutet, dass \verb|erhoeheMaybe| als Spezialfall
  die Funktion \verb|increaseByOne| aus \cref{fig:curryingExample} ist.
  Wir können dies in \cref{fig:maybeTypeAnwenden4} sehen.

  \begin{figure}[h]
    \begin{minted}[frame=single]{haskell}
    erhoeheMaybe [1, 2, 3] = [2, 3, 4]
    erhoeheMaybe Just 1 = Just 2
    \end{minted}
    \caption{Anwendung von \glqq erhoeheMaybe\grqq  auf Listen von Zahlen und auf \glqq Maybe Int\grqq .}%
    \label{fig:maybeTypeAnwenden4}
  \end{figure}

  Wir können dieses Schema fortführen und beispielsweise einen Datentypen
  für Bäume definieren, wobei ein Knoten jeweils einen Wert vom Typ \verb|a|
  beinhalten kann. So könnten wir auch eine Funktorinstanz für diesen
  Datentypen definieren, wodurch die Funktion \verb|erhoeheMaybe| auch auf
  Bäume angewendet werden kann. Wir sehen also, dass die Funktorialität
  von Datentypen ein sehr mächtiges Konzept ist.

  \subsubsection{Funktoren von \textbf{Hask} in andere Kategorien}

  Wichtig ist, dass wir bisher nur Funktoren von \textbf{Hask} in \textbf{Hask}
  gesehen haben. Wir nennen solche Funktoren \glqq Endofunktoren\grqq .
  Interessant ist, dass wir jedoch die Ähnlichkeiten zwischen Haskell
  und Rust auch durch Funktoren beschreiben können. Wir können allgemeiner
  \glqq Ähnlichkeiten\grqq  zwischen Haskell und anderen Programmiersprachen
  mit den mathematischen Konzepten, welche das Fundament von Haskell
  bilden, kurz und bündig beschreiben.

  Wir nehmen dazu an, dass wir analog zu der Kategorie \textbf{Hask}
  eine Kategorie \textbf{Rust} haben, welche analog zu \textbf{Hask}
  definiert ist. Die Objekte der Kategorie \textbf{Rust} sind die
  Typen und die Morphismen sind die Funktionen zwischen den Typen in
  Rust. Im allgemeinen würde dies keine Kategorie im Sinne der Definition
  von Kategorien aus der Einleitung sein, weshalb wir analog zu der
  Definition von \textbf{Hask} beispielsweise nur terminierende Funktionen
  in \textbf{Rust} zulassen. Die Konzepte, welche wir hier erschließen,
  lassen sich jedoch auch auf nicht-terminierende Funktionen erweitern,
  wobei wir darauf nicht weiter eingehen möchten.

  Wir haben gesehen, dass für Typen, welche eine Funktor-Instanz sind
  gilt, dass sie die Funktion \verb|fmap :: (a -> b) -> f a -> f b| haben,
  welche eine Funktion vom Typ \verb|a -> b| auf einen Funktion vom Typ
  \verb|f a -> f b| bis auf Currying schickt. Das wichtige dabei ist,
  dass wir dies \glqq strukturerhaltend\grqq  machen möchten. Betrachten wir
  dazu das Beispiel aus \cref{fig:maybeTypeAnwenden2}. Wir möchten,
  dass wenn \verb|x| ein Wert vom Typ \verb|Int| ist, dann soll
  zuerst das Anwenden von \verb|(+ 1)| auf \verb|x| und dann das
  Anwenden von \verb|Just| auf das Ergebnis das selbe sein, wie wenn
  man zuerst \verb|Just| auf \verb|x| anwendet und dann \verb|fmap (+ 1)|
  auf das Ergebnis. Wir können dies in \cref{fig:kommutativitaetFunktor}
  sehen.

  \begin{figure}[h]

    \begin{minted}[frame=single]{haskell}
    fmap (+ 1) (Just 1) = Just 2
    Just ((+ 1) 1) = Just 2
    \end{minted}
    \caption{Beispiel der grundlegenden Rechenregel für die Funktor-Instanz von \glqq Maybe\grqq .}%
    \label{fig:kommutativitaetFunktor}
  \end{figure}

  Eine andere Art diese ausschlaggebende Eigenschaft von Funktor-Instanzen
  zu beschreiben, ist indem wir ein sogenanntes \glqq kommutatives Diagramm\grqq 
  aufstellen. Wir können dies in \cref{fig:kommutativesDiagramm} sehen.
  Was wir damit aussagen möchten ist, wenn ich zuerst rechts und dann runter
  gehe, indem ich die entsprechenden Funktionen anwende, ist es das selbe,
  wie wenn ich zuerst runter und dann rechts gehe.

  \begin{figure}[h]
    % https://tikzcd.yichuanshen.de/#N4Igdg9gJgpgziAXAbVABwnAlgFyxMJZABgBoBGAXVJADcBDAGwFcYkQAdD2mAJwCMAPgDMABPUEgAvqXSZc+QijLFqdJq3ZceAwROmyQGbHgJFypVTQYs2iTtz5ChBuScXmKam5vvanIqIuUmowUADm8ESgwrwQALZIFiA4EEgATDQAFjD0UOyQYGwyMXGJiMmpSGQgOXlIYMyMjNYadg46QsKSJSCxCdU0VYgAzK22Wo66wvH0aKLhPYb95ZVpiOm9K0lD65nqE35TQouufWUZu0hjIIz0-DCMAArypkq3MMI4IOO+HQHdaSUKRAA
    \[
      \begin{tikzcd}
        \verb|a| \arrow[r, "\text{g}"] \arrow[d, "\text{f}"] & \verb|b| \arrow[d, "\text{f}"] \\
        \verb|f a| \arrow[r, "\text{fmap(g)}"]                   & \verb|f b|
      \end{tikzcd}
    \]
    \caption{kommutatives Diagramm, welches die Eigenschaften von \glqq fmap\grqq  darstellt.}%
    \label{fig:kommutativesDiagramm}
  \end{figure}
  
  Für Funktoren in \textbf{Hask} ist dies die einzige Eigenschaft, welche
  sie erfüllen müssen. Allgemein ist dies sogar die kategorientheoretische
  Definition eines Funktors. Ein Funktor $F$ von einer Kategorie $\mathcal{C}$ in
  eine Kategorie $\mathcal{D}$ ordnet jedem Objekt $a \in \mathcal{C}$ ein
  Objekt $F(a) \in \mathcal{D}$ zu und jedem Morphismus $f : a \to b$ in
  $\mathcal{C}$ einen Morphismus $F(f) : F(a) \to F(b)$ in $\mathcal{D}$ zu,
  sodass das folgende Diagramm kommutiert:

  \[
    \begin{tikzcd}
      a \arrow[r, "f"] \arrow[d, "F"] & b \arrow[d, "F"] \\
      F(a) \arrow[r, "F(f)"]                   & F(b)
    \end{tikzcd}
  \]

  Besitzt ein Typ in \textbf{Hask} eine Funktor-Instanz, 
  so sagen wir damit aus, dass die konstruieren dieses Typens Funktoren
  im Sinne der Kategorientheorie sind

  Dies erlaubt es uns nun auch einen abstrakteren Funktor $F$ von \textbf{Hask}
  nach \textbf{Rust} zu definieren, indem wir jeden Typen $a$ in \textbf{Hask}
  mit einem Typen $F(a)$ in \textbf{Rust} assoziieren. Beispielhaft könnte
  der Typ \verb|Int| mit dem Typen \verb|u32| assoziiert werden und so weiter.
  Insbesondere müsste dann auch $F(\verb|Maybe Int|)$ mit einem Typen
  in \textbf{Rust} assoziiert werden, welcher die Funktionalität von
  \verb|Just| erhält. Dafür bietet sich der Typ \verb|Option<u32>| in
  \textbf{Rust} an. Der Typ \verb|Option<u32>| in Rust ist dadurch
  definiert, dass er entweder den Wert \verb|None| oder den Wert
  \verb|Some(u32)| annehmen kann. Wir können somit $F(\verb|Just x|)$
  auf \verb|Some(x)| für ein beliebiges $x$ von Typ \verb|Int| in \textbf{Hask}
  abbilden, sowie $F(\verb|Nothing|)$ auf \verb|None| abbilden.
  Nun müssen wir nur noch für eine Abbildung \verb|g :: Int -> b| in \textbf{Hask}
  die zugehörige Abbildung \verb|F(g)| in \textbf{Rust} finden.
  Dies ist jedoch dadurch klar, indem wir durch Pattern-Matching auf
  \verb|Some(x)| und \verb|None| unterscheiden und für beide Fälle
  die entsprechende Abbildung \verb|g| anwenden. Dies ist in \cref{fig:OptionFunktor}
  dargestellt.

  \begin{figure}[h]
    \begin{minted}[frame=single]{haskell}
    // result ist ein Option<u32>
    match result {
      Some(x) => F(h(x)),
      None    => None
      }
    \end{minted}
    \caption{Implementation von $F(\text{g} :: \text{Maybe Int} \to \text{a})$, wobei h die eindeutig von f induzierte Funktion $\text{h} :: \text{Just Int} \to a$ ist.}%
    \label{fig:OptionFunktor}
  \end{figure}

  Was dies bedeutet ist, dass wir mit den Konzepten, welche in Haskell
  existieren vergleiche mit anderen Programmiersprachen ziehen können.
  Insbesondere können wir sehen, das \textbf{Rust} Strukturen
  besitzt, welche denen in \textbf{Hask} ähneln. Auch können wir sehen,
  dass es in Haskell einfacher möglich ist mit verkapselten Typen zu
  arbeiten, da wir auf dem \verb|Maybe| Typen eine Funktor-Instanz
  besitzen. Diese fehlt jedoch in Rust, wodurch wir nur durch Pattern-Matching
  ein ähnliches Verhalten auf den \verb|Option| Typen erreichen können.
  Hier zeigt sich jedoch ein Muster, welches Haskell vollkommen durchzieht:
  Konzepte, welche in Haskell zum eleganten Programmieren existieren
  haben auch praktische Anwendung in anderen Bereichen, wie zum Beispiel
  dem Vergleichen von Programmiersprachen. Haskell zu lernen trägt somit
  dazu bei, generell besser Strukturen zu erkennen, was im Endeffekt
  auch zu einem besseren Programmieren führt.

  \section{Monaden und unser erstes Haskell Programm}

  \epigraph{
    If you're approaching Haskell monads a little nervously, that's understandable. You've probably heard that monads are a very powerful code-structuring technique, but ... sometimes, with great power comes great (apparent) nonsensicality.
  }{\textit{Haskell Wiki \\ All About Monads}}

  Ein sehr Populäres Konzept in Haskell, woran viele Programmierer beim
  Erlernen der Sprache scheitern, sind Monaden. Monaden sind ein
  weiteres Konzept aus der Kategorientheorie, welches in Haskell
  allgegenwärtig ist. Fürwahr, wir sind nicht in der Lage ein laufendes
  Programm in Haskell zu schreiben, ohne ein zumindest grundlegendes
  Verständnis von Monaden vorzuweisen.
  Monaden in Haskell sind ein weiterer Datentyp
  $\verb|type Monad :: (* -> *) -> Constraint|$, welche eine Funktor-Instanz
  besitzen, also auch Funktoren sind, und
  zusätzlich eine Operation $\verb|return :: a -> m a|$ und eine
  Operation \verb|>>= :: m a -> (a -> m b) -> m b|
  besitzen. Die Operation \verb|return| ist eine sogenannte
  \glqq natürliche Transformation\glqq , welche ein Objekt in eine Monade
  steckt. Wir werden nicht weiter darauf eingehen, was eine natürliche
  Transformation ist, jedoch ist es nützlich zu wissen, dass diese
  genau so wie Funktoren als eine Art von Abbildung zwischen Kategorien
  (in unserem Fall von \textbf{Hask} nach \textbf{Hask}) definiert
  werden können, welche gewisse Eigenschaften besitzen.
  Die Operation \verb|>>=| ist eine sogenannte \glqq Bindung\glqq , welche
  es uns erlaubt \glqq gefangene Werte\grqq  in einer Monade zu manipulieren.

  Betrachten wir jedoch ein paar Beispiele, um uns diesem gefürchteten
  Konzept zu nähern und zu sehen, dass es gar nicht so angsteinflößend
  ist.

  \subsection{Listen als Monade}

  Wir haben bereits ein inniges Verständnis von Listen in Haskell
  in den vorherigen Abschnitten aufgebaut. Wir wissen, dass Listen
  eine Funktor-Instanz besitzen, welche uns erlaubt Listen zu
  manipulieren, indem wir Funktionen für die Unterliegenden Datentypen
  in der Liste definieren.

  Wir möchten nun einen Schritt weiter gehen und uns anschauen, warum
  Listen auch eine kanonische Monaden-Instanz in Haskell besitzen,
  also eine Instanz, welche in der Standardbibliothek \verb|Prelude|
  von Haskell definiert ist. Außerdem stellt sich die Frage, warum
  wir überhaupt eine Monaden-Instanz für Listen brauchen.

  Um Listen zu einer Monade zu machen, müssen wir uns zuerst genauer ansehen,
  was eine Monade ist. Eine Monade ist ein Datentyp, dessen Definition
  in \cref{fig:MonadeDef} dargestellt ist.

  \begin{figure}[h]
    \begin{minted}[frame=single]{haskell}
    type Monad :: (* -> *) -> Constraint
    class Applicative m => Monad m where
      (>>=) :: m a -> (a -> m b) -> m b
      (>>) :: m a -> m b -> m b
      return :: a -> m a
      {-# MINIMAL (>>=) #-}
      -- Defined in 'GHC.Base'
    \end{minted}
    \caption{Definition der \glqq Monaden\grqq  Typenklasse.}%
    \label{fig:MonadeDef}
  \end{figure}

  Die Definition einer Monade ist hierbei genau so wie die in
  \cref{fig:functorTypeClass} aufgeschrieben und sollte auch so
  verstanden werden. Wir haben hier jedoch eine Einschränkung,
  welche dadurch gegeben ist, dass hier \verb|Applicative m => Monad m|
  steht. Dies stellt einen Type-Constraint da, welche
  und bereits bekannt aus dem Abschnitt über Funktoren sind.
  Dieser Type-Constraint besagt, dass ein Typ, welcher eine Monade
  ist, nur dann eine Monade sein kann, wenn dieser auch ein \glqq Applicative\grqq  ist.
  Die Definition eines Applicatives ist in \cref{fig:ApplicativeDef}
  auffindbar. Relevant ist hierbei nur, dass auch ein Applicative einen
  Type-Constraint besitzt, welcher besagt, dass jeder Typ, welcher
  eine Instanz von Applicative ist, auch eine Instanz von Functor
  sein muss. Die Operationen \verb|<*>|, \verb|GHC.Base.liftA2|,  \verb|*>| 
  und \verb|<*| sollten nach unserer Analyse der Funktor Typenklasse bereits
  verständlich sein, wobei ihr Verständnis keine Anforderung
  für das Verständnis von Monaden für den Rest dieses Abschnittes ist.

  \begin{figure}[h]
    \begin{minted}[frame=single]{haskell}
    class Functor f => Applicative f where
      pure :: a -> f a
      (<*>) :: f (a -> b) -> f a -> f b
      GHC.Base.liftA2 :: (a -> b -> c) -> f a -> f b -> f c
      (*>) :: f a -> f b -> f b
      (<*) :: f a -> f b -> f a
      {-# MINIMAL pure, ((<*>) | liftA2) #-}
        -- Defined in 'GHC.Base'
    \end{minted}
    \caption{Definition der \glqq Applicative\grqq  Typenklasse.}%
    \label{fig:ApplicativeDef}
  \end{figure}

  Wir nehmen nun an, dass Listen eine Instanz von Applicative sind.
  Dies ist auch der Fall, jedoch werden wir dies nicht beweisen. Der
  Beweis, dass Listen eine Instanz von Monaden sind wird dadurch geführt,
  indem wir konkrete Implementierungen für die Operationen
  \verb|>>=| und \verb|return| definieren.

  Betrachten wir zunächst die Operation \verb|return|, welche einen
  Wert in eine Monade steckt. Was dies bedeutet, ist bei Listen relativ
  selbsterklärend. Sei nun \verb|x| ein Wert vom Typ \verb|a|, dann
  ist \verb|return x| eine Liste, welche nur den Wert \verb|x| enthält.
  Dies stimmt auch mit der Typen-Signatur von \verb|return| überein,
  welche \verb|return :: a -> m a| ist.

  Nun müssen wir nur noch die Operation \verb|>>=| definieren. Diese
  Operation soll uns erlauben, einen Wert aus einer Monade zu ziehen,
  dann zu bearbeiten und dann zurück in die Monade zu legen.
  Sei nun \verb|xs| eine Liste vom Typ \verb|[a]| und \verb|f| eine
  Funktion vom Typ \verb|a -> b|. So können wir \verb|(>>=) xs f = fmap f xs|,
  beziehungsweise kompakter \verb|>>= = fmap| definieren. Dies ist
  uns bereits aus dem Abschnitt über Funktoren bekannt.

  Wir sehen hiermit, dass Monaden es uns erlauben eine Art \glqq Gefängnis\grqq 
  für unsere Werte zu schaffen, in welchem wir die Werte manipulieren
  können. Jedoch gibt es keine Operation mit der Typen-Signatur
  \verb|get :: m a -> a|, welche uns erlaubt, einen Wert aus der
  Monade zu ziehen. Dies ist auch bei allgemeinen Monaden nicht möglich,
  wobei wir eine solche Funktion bei Listen relativ leicht wie in
  \cref{lst:listeget} definieren können. Wichtig ist in diesem Beispiel,
  dass die Funktion nicht total ist, also nur für den Fall einer einelementigen
  Liste definiert ist. Sollte man eine Liste mit mehr oder weniger als
  einem Element übergeben, so wird eine Laufzeitfehler auftreten.

  \begin{figure}[h]
    \begin{minted}[frame=single]{haskell}
    get :: [a] -> a
    get [x] = x
    \end{minted}
    \caption{Eine Funktion, welche einen Wert aus einer Liste und somit aus einem Objekt, welches eine Monadenstruktur besitzt, herausnimmt.}%
    \label{fig:listget}
  \end{figure}

  \subsection{Die IO Monade und Hello World in Haskell}

  Die wohl wichtigste Monade ist die \verb|IO| Monade. Diese Monade
  erlaubt es uns, Eingaben und Ausgaben zu machen, woher auch der
  Name \verb|IO| kommt. Ein wichtiges Beispiele für die Verwendung der
  \verb|IO| Monade ist in \cref{fig:helloworld} zu finden. Hier definieren
  wir den Wert \verb|main :: IO ()|, welcher der Start eines jeden
  Haskell Programms ist. Der Klarheit halber sei erwähnt, dass \verb|IO|
  also mit dem Typen \verb|()| als Parameter verwendet wird, welcher
  den Typen \verb|Unit| in anderen Sprachen entspricht und lediglich
  ein leeres Tupel der Länge 0 darstellt.

  \begin{figure}[h]
    \begin{minted}[frame=single]{haskell}
    main :: IO ()
    main = putStrLn "Hello World"
    \end{minted}
    \caption{Hello World in Haskell}%
    \label{fig:helloworld}
  \end{figure}

  Zeigen wir eine Parallelität zwischen der \verb|IO| Monade und der
  \verb|[]| Monade, der Monade der Listen, auf um die Funktionalität
  der \verb|IO| Monade zu verstehen. Wir können uns vorstellen, dass
  \verb|IO| eine Art \glqq Gefängnis\grqq  für Werte ist, in diesem Fall 
  von \verb|IO ()| für leere Tupel, genau so wie \verb|[a]| eine Art
  \glqq Gefängnis\grqq  für Werte vom Typ \verb|a| ist. Allerdings gibt es einen
  bedeutenden Unterschied zwischen den beiden Monaden. Während wir
  in \verb|[a]| wissen, was genau in unserem \glqq Gefängnis\grqq  ist, so
  wissen wir dies bei \verb|IO| nicht. In \verb|[Int]| befindet sich
  für ein Element \verb|[1,2,3]| beispielsweise nur die Zahl 1, die Zahl
  2 und die Zahl 3 in dem Gefängnis. In dem Typen \verb|IO ()|
  gibt es konstruktionsgemäß, weil der Typ \verb|()| nur das Element \verb|()|
  besitzt, nur das Element \verb|IO ()|. Dieses Element kann jedoch
  neben sich in seinem \glqq \verb|IO| Gefängnis\grqq  noch weitere Sachen sitzen haben.
  Das einzige was unser \glqq \verb|IO| Gefängnis\grqq  garantieren muss ist, dass
  der Binding-Operator \verb|>>=| und der \verb|return| Operator
  ihre Rechenregeln beibehalten müssen. An dem Beispiel von Hello World
  in \cref{fig:helloworld} können wir uns dies gut vorstellen. Lässt man
  dieses Programm laufen, so wird die Ausgabe \verb|Hello World| auf
  dem Bildschirm erscheinen. Wir haben somit einen Seiteneffekt beim
  ausführend unseres Programmes, also beim \glqq Aufrufen des Wertes \verb|main|\grqq  erzeugt.

  Nun gibt es verschiedene philosophische Herangehensweisen, wie wir
  dies Interpretieren können. Wir können uns beispielsweise vorstellen,
  dass in unserem \glqq \verb|IO| Gefängnis\grqq  der Wert \verb|()| nun zusammen mit
  dem Seiteneffekt, welcher \glqq \verb|Hello World|\grqq  auf dem
  Bildschirm ausgibt, sitzt.
  Dies bedeutet, wenn der Wert \verb|main| aufgerufen wird, so wird
  das gesamte \glqq \verb|IO| Gefängnis\grqq  betrachtet, was den Seiteneffekt
  anzeigt.

  Das wichtige hierbei ist jedoch, dass wir nun eine rigorose Theorie
  für Seiteneffekte haben, welche wir in Haskell nun nutzen können
  um mit Seiteneffekten, wie Beispielsweise der Ausgabe auf dem Bildschirm,
  umzugehen.

  Möchte man Beispielsweise mehrere Seiteneffekte in seinem Programm
  haben, so können wir den uns bekannten Binding-Operator \verb|>>=|
  verwenden. In \cref{fig:ioexamples} sehen wir ein Beispiel, welches zwei
  Seiteneffekte in den Wert \verb|main| einbindet. Lässt man dieses
  Programm nun laufen, so wird zuerst \verb|Hello| auf dem Bildschirm
  ausgegeben und anschließend \verb|World|.
  
  \begin{figure}[h]
    \begin{minted}[frame=single]{haskell}
    main :: IO ()
    main = print "Hello" >>= \x -> print "World"
    \end{minted}
    \caption{Hello World in Haskell, indem wir die Monadenstruktur durch den Binding-Operator ausnutzen.}%
    \label{fig:ioexamples}
  \end{figure}

  Zuallerletzt können wir auch den \verb|return| Operator nutzen, um
  unser Programm nichts machen zu lassen. In \cref{fig:ioexamples2} sehen
  wir ein Beispiel dafür. Das Kontatinieren von Seiteneffekten 
  zu neuen Seiteneffekten und die Existenz eines Seiteneffekts,
  welcher nichts macht, also der \glqq nichts-machende Seiteneffekt\glqq  \verb|return ()|
  vom Typen \verb|IO ()|, sind die Eigenschaften, welche
  intuitiv für Seiteneffekte klar sind und Seiteneffekte ausmachen.
  Und diese sind genau die Eigenschaften, welche wir mit der Monadenstruktur
  erfassen.

  Eine interessante Folgerung, welche aus dem vorherigen Paragraphen
  klar erkenntlich ist, ist hierbei dass Monaden ein Monoidobjekt
  in der Kategorie der Endofunktoren sind. Wir werden diesen Satz nicht
  beweisen, jedoch in den Zusammenhang mit der Kategorie \textbf{Hask}
  und der Monade \verb|IO| bringen.
  Ein Objekt besitzt eine Monoidstruktur,
  wenn man je zwei Elemente daraus zu einem dritten Element daraus
  verknüpfen kann.
  Für die Monade \verb|IO| übernimmt diese Rolle der Verknüpfung der Binding-Operator.
  Dieser muss zudem noch assoziativ sein, was intuitiv für den Binding-Operator
  klar ist, da es das selbe ist, wenn $a, b, c$ drei Seiteneffekte sind,
  ob man zuerst $a$ gefolgt von $b$ mit $c$ ausführt oder zuerst $a$ mit $b$
  gefolgt von $c$ ausführt. Die letzte Eigenschaft welche für eine Monoidstruktur
  benötigt wird ist die Existenz eines neutralen Elements, welches
  also nichts macht. Dieses haben wir bereits im vorherigen Paragraphen
  konstruiert, nämlich den Wert \verb|return ()| vom Typen \verb|IO ()|.

  \subsection{Warum Monaden? Oder: Warum Haskell?}

  Schlussendlich stellt sich jedoch die Frage, warum wir einen solchen 
  Aufwand betreiben, um beispielsweise mit Seiteneffekten umzugehen.
  Die Antwort ist jedoch ganz einfach. Wir haben am Anfang damit begonnen,
  einfache Typen zu erstellen. Stets waren wir bemüht, die Theorie, welche
  wir aufbauen rigoros und wohldefiniert zu halten. Haben wir eine
  Struktur gefunden, so haben wir versucht dieser einen Namen zu geben und
  diese zu verallgemeinern, ganz im Sinne der \glqq Modularität\grqq  und
  \glqq Ausdrucksfähigkeit\grqq , so wie es in \cite{abelson1996StructureAndInterpretationOfComputerPrograms}
  empfohlen wird.
  Speziell haben wir dies bei der \verb|Functor|
  Typenklasse gesehen, welche das Konzept \glqq Abbildungen zwischen Strukturen
  können auf ähnliche Abbildungen auf ähnlichen Strukturen geschickt werden\glqq 
  darstellt. Nun war es am Ende eigentlich nur ein glücklicher
  Zufall, dass unsere Struktur der Monaden genau die Eigenschaften,
  welche Seiteneffekte ausmachen, besitzt.

  Die Bedeutung letztendlich ist, dass wir in Haskell genau wissen möchten,
  womit wir arbeiten.
  Während man sich in anderen Programmiersprachen
  wie Python oder Java damit begnügt, dass Seiteneffekte existieren
  und immer auftreten können,
  selbst wenn wir explizit angeben,
  dass unsere Funktion einen Wert von Typ $a$ einnimmt und einen Wert
  von Typ $b$ zurückgibt,
  kann es immer noch passieren, dass die Funktion plötzlich als Seiteneffekt
  etwas anderes macht.
  In Haskell wissen wir jedoch, dass wenn
  wir eine Typensignatur angeben, die Funktion auch genau diesen Typ
  hat\footnote{außer bei partiellen Funktionen, welche Definitionslückene aufweisen, oder nicht-terminierenden Funktionen}.

  Genau so, wie wir sagen könnten \glqq uns ist egal was ein 'Körper` ist,
  hauptsache wir wissen was 1,2,3, \dots sind\glqq , könnten wir auch sagen
  \glqq uns ist egal was eine Monade ist, hauptsache wir wissen dass es
  Seiteneffekte gibt\glqq . Wir können auch ohne das nötige Hintergrundwissen
  Programmieren, jedoch können wir deutlich eleganteren und präziseren Code
  schreiben, wenn wir die den Problemen grundlegenden Strukturen kennen.

  Haskell hilft uns dadurch besser zu programmieren, indem
  es uns dazu zwingt diese Konzepte zu verstehen und umzusetzen.
  Hier können wir nun auch sehen, warum Haskell zwar sehr eleganten
  und effizienten Code erzeugen kann, jedoch in der Industrie nicht
  weit verbreitet ist. Oft ist es einfacher und schneller, eine
  \glqq quick and dirty\grqq  Lösung zu finden welche \glqq einfach funktioniert\grqq ,
  als sich mit den Grundlagen seiner Probleme auseinanderzusetzen.
  Insbesondere ist die Hürde für die Einarbeitung in Haskell sehr hoch,
  was bei anderen populären Sprachen wie Python oder JavaScript nicht der Fall ist.

  Schlussendlich ist es für jeden Programmierer ratsam in seinem
  Abenteuer durch die Welt der Programmiersprachen auch einen Zwischenhalt
  bei Haskell einzulegen. Die Konzepte welche in Haskell behandelt
  werden, können den Programmierstil in anderen Sprachen, welche
  der Programmierer verwendet, positiv beeinflussen.
  Auch wenn der Programmierer nicht unbedingt Haskell am Ende in
  Production verwendet, so bieten die Ideen, welche Haskell vermittelt,
  allein schon genügend Anreiz zum erlernen der Sprache.

%   \subsection{Unterthemen}.
%   Unterthemen geben bereits eine Aufteilung in Abschnitte vor und jedes
%   Unterthema sollte nicht mehr als einen Abschnitt umfassen.
%   Dabei sollte nach dem Titel des Abschnitts je ein
%   kurzer, einführender Text stehen, bevor der Abschnitt in weitere,
%   sinnvolle Unterabschnitte aufgeteilt wird.
%   Dieser einführende Text steht direkt hinter dem jeweiligen Titel und
%   nicht in einem etwaigen Unterabschnitt.


  % \section{Abbildungen, Tabellen und Beispiele}

  % Sie werden in Ihren Zusammenfassungen nicht darum herum kommen,
  % Abbildungen, Tabellen und Beispiele zur besseren Anschaulichkeit zu nutzen.
  % Diese Abbildungen müssen im Text referenziert und besprochen werden
  % und sollten nicht alleine stehen.
  % Ebenfalls sollten Abbildungen und Tabellen zentriert dargestellt werden
  % (mit \texttt{\textbackslash centering}).

  % Mit \cref{fig:initial-draft} wird beispielhaft die Referenzierung einer Grafik
  % demonstriert.
  % Auf \LaTeX-Seite geschieht dies durch das Setzen eines Labels
  % (\texttt{\textbackslash label}) für die Grafik
  % und durch die Nutzung der Befehle
  % \texttt{\textbackslash ref} oder \texttt{\textbackslash cref}.

  % \begin{figure}[h]
  %   \centering
  %   \includegraphics[width=4cm]{fig/hhulogo.pdf}
  %   \caption{Das neue HHU-Logo.}%
  %   \label{fig:initial-draft}
  % \end{figure}

  % Der \texttt{\textbackslash cref}-Befehl entstammt hierbei dem
  % \texttt{cleveref}-Paket und setzt automatisch den korrekten Bezeichner
  % vor die Referenz, in diesem Beispiel \grqq Abbildung\glqq .
  % Zum Vergleich: \texttt{\textbackslash ref} setzt nur die
  % Nummer \ref{fig:initial-draft}, die der Abbildung zugeteilt wurde.

%   \subsection{Beispiele}
%   Viele Themen profitieren von einem guten Beispiel,
%   um Konzepte oder Eigenschaften klarer zu präsentieren.
%   Achten Sie bei Beispielen darauf,
%   dass diese nicht unnötig komplex sind und nur auf das nötige reduziert sind.
%   Sollte es sich anbieten, ist es hilfreich ein Running Example zu nutzen,
%   welches von allen Unterthemen aufgegriffen wird.
%   So ist der Leser bereits mit den Gegebenheiten des Beispiels vertraut
%   und kann Unterschiede zwischen den verschiedenen Szenarien genauer
%   differenzieren.


  % \subsection{Code}

  % \begin{figure}[h]

% \begin{lstlisting}
% public class HelloWorld 
% {
 
  %      public static void main (String[] args)
  %      {
  %            // Ausgabe Hello World!
  %            System.out.println("Hello World!");
  %      }
% }
% \end{lstlisting}
 % \caption{So binden Sie Code ein}%
  % \label{fig:codeBlock}
% \end{figure}

  % \subsection{Tabellen}

% \begin{table}[h!]
% \centering
% \begin{tabular}{ |c|c|c| } 
 % \hline
 % Sprache & Jahr & Version \\ 
% \hline
 % Java & 1997 & JSE 19.0 \\ 
 % C & 1972 & C17 \\ 
 % \hline
% \end{tabular}
% \caption{Eine sinnfreie Tabelle}
% \label{table:1}
% \end{table}

  % \section{Literaturverweise}

  % Die Angabe von Quellen und Verweisliteratur ist essentiell in Ihrer
  % Ausarbeitung,
  % als auch später in Ihrer Bachelorarbeit.

  % Literaturverweise werden in \LaTeX{} mittels
  % \texttt{\textbackslash cite} getätigt.
  % Verweise stehen hinter der Aussage, die sie stützen sollen,
  % oder hinter dem Namen eines referenzierten Konzepts oder Eigenschaft
  % und gehören in den Satz hinein.
  % Beachten Sie folgendes Beispiel:
  % \begin{itemize}
  %   \item In \grqq Artificial Intelligence: A Modern Approach\glqq ~\cite{russell2002artificial}
  %     werden moderne Ansätze und Techniken
  %     der künstlichen Intelligenz vorgestellt.
  %   \item bei direkten Zitationen wir zusätzlich die Seite angegeben wo das Zitat zu finden ist, z.B. \grqq ...\glqq ~\cite[p.1]{russell2002artificial}
  % \end{itemize}

  % Am Ende der Ausarbeitung werden die referenzierten Quellen automatisch
  % aufgelistet.

  % \subsection{Bibtex-Einträge und \texttt{references.bib}}

  % Um mit \texttt{\textbackslash cite} referenzieren zu können,
  % müssen die entsprechenden Einträge in der Datei \texttt{references.bib}
  % stehen.
  % Sie finden in der Vorlage bereits den beispielhaften Eintrag für
  % \grqq Artificial Intelligence: A Modern Approach\glqq ~\cite{russell2002artificial}.

  % Für Literaturquellen empfiehlt es sich in aller Regel
  % \url{https://scholar.google.com} zu nutzen.
  % Hier kann man sich direkt den Bibtex-Eintrag einer Quelle
  % angeben lassen und muss diesen so nicht selbst zusammenstellen.

  % Beachten Sie auch den Artikel
  % \grqq Common Errors in Bibliographies\glqq 
  % von John Owens.\footnote{%
  %   \url{https://www.ece.ucdavis.edu/~jowens/biberrors.html}}
  % Des Weiteren ist der Vortrag von Simon Peyton Jones \href{https://www.microsoft.com/en-us/research/academic-program/write-great-research-paper/}{\grqq How to write a great research paper\glqq } zu empfehlen.

  \appendix
  % \section{Beispielhafter Anhang}
  % Der Anhang wird mit dem \texttt{\textbackslash{}appendix}-Befehl eingeleitet.
  % Alle folgenden Sektionen werden als Anhang formatiert.

\backmatter
\listoffigures
% \listoftables
\bibliography{references}
\bibliographystyle{plaindin}

\end{document}

% Quellen

% Why You’ll Probably Never Use Haskell in Production
% https://henrydangprg.com/2019/05/27/why-youll-probably-never-use-haskell-in-production/

% Choosing Haskell isn’t a stand-in for good software design
% https://ozataman.medium.com/choosing-haskell-isnt-a-stand-in-for-good-software-design-7d893882f963

% A History of Haskell: Being Lazy With Class
% https://dl.acm.org/doi/10.1145/1238844.1238856

% Funktionale Programmierung mit Haskell/ Typklassen
% https://de.wikibooks.org/wiki/Funktionale_Programmierung_mit_Haskell/_Typklassen

% Making Our Own Types
% https://mmhaskell.com/liftoff/data-types

% Functor
% https://wiki.haskell.org/Functor

% Kind
% https://wiki.haskell.org/Kind

% Type constructor
% https://wiki.haskell.org/Constructor

% Currying
% https://de.wikipedia.org/wiki/Currying
